% GENERATED FILE. Edit canonical lessons, then run npm run generate:books.
% canonical-source-hash: fnv1a64:82e3ecc9dc72ad61
% canonical-lessons: IT-C448-panchina, IT-C448-lampione, IT-C448-cestino, IT-C448-tombino, IT-C448-ingombrare

\chapter{Panchine e lampioni}
\label{ch:it-panchine-e-lampioni}
\hlchaptermodality{448}

\begin{chapteropening}
\textbf{By the end of this chapter:} \emph{I can describe things on the street: a bench, a street lamp, a litter bin, a drain cover and things in the way.}

Complete the last lesson of chapter 448: I can describe things on the street: a bench, a street lamp, a litter bin, a drain cover and things in the way.
\end{chapteropening}

\section[la panchina]{la panchina — the bench (in a park); the substitutes' bench}

\label{lesson:IT-C448-panchina}



\begin{quote}
Before the new one: say the Italian for to report, then the Italian for the identity card.
\end{quote}

\subsection*{You'll want to know: la panchina}
\textbf{la panchina} — \textquotedblleft{}the bench (in a park); the substitutes' bench\textquotedblright{}.

\begin{quote}
Ci sediamo su una panchina al sole.
\end{quote}

\begin{quote}
\emph{We sit on a bench in the sun.}
\end{quote}

In football, \emph{in panchina} — on the bench.

\begin{cousinweb}
\emph{panchina} is a little \emph{panca} 'bench', which was borrowed from Lombardic \emph{panch} 'bench' (EtymDB), a Germanic word.
\end{cousinweb}

\begin{grammarlens}[title={What you've built}]
One more word for telling, explaining and describing at B1.
\end{grammarlens}

\subsection*{Your turn}
\begin{itemize}
\raggedright
  \item \emph{Say it:} \emph{la panchina}
  \item \emph{Say it:} \emph{la panchina}, once more
  \item \emph{Say it:} a sentence of your own with \emph{la panchina}
  \item \emph{Recall:} say the Italian for to report, then the Italian for the identity card, then say \emph{la panchina} again
\end{itemize}

\begin{tcolorbox}[breakable,colback=teal!4,colframe=teal!35!black,title={Before you move on}]
What does la panchina mean? (\textquotedblleft{}The bench (in a park); the substitutes' bench\textquotedblright{}.) Say it once more.
\end{tcolorbox}
\section[il lampione]{il lampione — the street lamp}

\label{lesson:IT-C448-lampione}



\begin{quote}
Before the new one: say the Italian for the identity card, then the Italian for the bench.
\end{quote}

\subsection*{You'll want to know: il lampione}
\textbf{il lampione} — \textquotedblleft{}the street lamp\textquotedblright{}.

\begin{quote}
Il lampione davanti a casa è spento.
\end{quote}

\begin{quote}
\emph{The street lamp in front of the house is off.}
\end{quote}

A big \emph{lampa}, a lamp.

\begin{cousinweb}
\emph{lampione} is a big \emph{lampa} 'lamp' (EtymDB), which was borrowed from Old French \emph{lampe}, from Latin \emph{lampas} and Greek \emph{lampás} 'torch'.
\end{cousinweb}

\begin{grammarlens}[title={What you've built}]
One more word for telling, explaining and describing at B1.
\end{grammarlens}

\subsection*{Your turn}
\begin{itemize}
\raggedright
  \item \emph{Say it:} \emph{il lampione}
  \item \emph{Say it:} \emph{il lampione}, once more
  \item \emph{Say it:} a sentence of your own with \emph{il lampione}
  \item \emph{Recall:} say the Italian for the identity card, then the Italian for the bench, then say \emph{il lampione} again
\end{itemize}

\begin{tcolorbox}[breakable,colback=teal!4,colframe=teal!35!black,title={Before you move on}]
What does il lampione mean? (\textquotedblleft{}The street lamp\textquotedblright{}.) Say it once more.
\end{tcolorbox}
\section[il cestino]{il cestino — the litter bin; the small basket}

\label{lesson:IT-C448-cestino}



\begin{quote}
Before the new one: say the Italian for the bench, then the Italian for the street lamp.
\end{quote}

\subsection*{You'll want to know: il cestino}
\textbf{il cestino} — \textquotedblleft{}the litter bin; the small basket\textquotedblright{}.

\begin{quote}
Butta la carta nel cestino.
\end{quote}

\begin{quote}
\emph{Throw the paper in the bin.}
\end{quote}

\emph{Il cestino della spesa} — a shopping basket.

\begin{cousinweb}
\emph{cestino} is a little \emph{cesta} 'basket', which is inherited from Latin \emph{cista}, from Greek \emph{kístē} 'box, basket' (EtymDB).
\end{cousinweb}

\begin{grammarlens}[title={What you've built}]
One more word for telling, explaining and describing at B1.
\end{grammarlens}

\subsection*{Your turn}
\begin{itemize}
\raggedright
  \item \emph{Say it:} \emph{il cestino}
  \item \emph{Say it:} \emph{il cestino}, once more
  \item \emph{Say it:} a sentence of your own with \emph{il cestino}
  \item \emph{Recall:} say the Italian for the bench, then the Italian for the street lamp, then say \emph{il cestino} again
\end{itemize}

\begin{tcolorbox}[breakable,colback=teal!4,colframe=teal!35!black,title={Before you move on}]
What does il cestino mean? (\textquotedblleft{}The litter bin; the small basket\textquotedblright{}.) Say it once more.
\end{tcolorbox}
\section[il tombino]{il tombino — the drain cover, the manhole}

\label{lesson:IT-C448-tombino}



\begin{quote}
Before the new one: say the Italian for the street lamp, then the Italian for the litter bin.
\end{quote}

\subsection*{You'll want to know: il tombino}
\textbf{il tombino} — \textquotedblleft{}the drain cover, the manhole\textquotedblright{}.

\begin{quote}
Con la pioggia i tombini sono pieni.
\end{quote}

\begin{quote}
\emph{With the rain the drains are full.}
\end{quote}

A little \emph{tomba}: a covered hole in the street.

\begin{cousinweb}
\emph{tombino} is a little \emph{tomba} 'tomb', which is inherited from Late Latin \emph{tumba}, from Greek \emph{týmbos} 'burial mound' (EtymDB).
\end{cousinweb}

\begin{grammarlens}[title={What you've built}]
One more word for telling, explaining and describing at B1.
\end{grammarlens}

\subsection*{Your turn}
\begin{itemize}
\raggedright
  \item \emph{Say it:} \emph{il tombino}
  \item \emph{Say it:} \emph{il tombino}, once more
  \item \emph{Say it:} a sentence of your own with \emph{il tombino}
  \item \emph{Recall:} say the Italian for the street lamp, then the Italian for the litter bin, then say \emph{il tombino} again
\end{itemize}

\begin{tcolorbox}[breakable,colback=teal!4,colframe=teal!35!black,title={Before you move on}]
What does il tombino mean? (\textquotedblleft{}The drain cover, the manhole\textquotedblright{}.) Say it once more.
\end{tcolorbox}
\section[ingombrare]{ingombrare — to clutter, to block (a space)}

\label{lesson:IT-C448-ingombrare}



\begin{quote}
Before the new one: say the Italian for the litter bin, then the Italian for the drain cover.
\end{quote}

\subsection*{You'll want to know: ingombrare}
\textbf{ingombrare} — \textquotedblleft{}to clutter, to block (a space)\textquotedblright{}.

\begin{quote}
Le biciclette ingombrano il marciapiede.
\end{quote}

\begin{quote}
\emph{The bikes are blocking the pavement.}
\end{quote}

\emph{Ingombrante} — bulky.

\begin{cousinweb}
\emph{ingombrare} is built from \emph{in-} and an old word for a heap or an obstacle; the origin of that word is not settled.
\end{cousinweb}

\begin{grammarlens}[title={What you've built}]
One more word for telling, explaining and describing at B1.
\end{grammarlens}

\subsection*{Your turn}
\begin{itemize}
\raggedright
  \item \emph{Say it:} \emph{ingombrare}
  \item \emph{Say it:} \emph{ingombrare}, once more
  \item \emph{Say it:} a sentence of your own with \emph{ingombrare}
  \item \emph{Recall:} say the Italian for the litter bin, then the Italian for the drain cover, then say \emph{ingombrare} again
\end{itemize}

\begin{tcolorbox}[breakable,colback=teal!4,colframe=teal!35!black,title={Before you move on}]
What does ingombrare mean? (\textquotedblleft{}To clutter, to block (a space)\textquotedblright{}.) Say it once more.
\end{tcolorbox}
